\documentclass{article}

% Language setting
\usepackage[german]{babel}
\usepackage{float}    
% Set page size and margins
% Replace `letterpaper' with `a4paper' for UK/EU standard size
\usepackage[letterpaper,top=2cm,bottom=2cm,left=3cm,right=3cm,marginparwidth=1.75cm]{geometry}
\usepackage{multicol}
\setlength{\columnsep}{1cm}
\usepackage{xcolor}
\usepackage{graphicx}
\usepackage{tcolorbox}
\tcbuselibrary{skins}
\usepackage{lipsum}

\definecolor{boxTitle}{HTML}{B1CF5F}
\definecolor{rose}{HTML}{CC7E85}
\definecolor{rose-d}{HTML}{83343C}
\definecolor{rose-l}{HTML}{DCA7AC}
\definecolor{boxBackground}{HTML}{DEF4C6}
\definecolor{boxFrame}{HTML}{004346}

\tcbset{my box/.style={
    enhanced, fonttitle=\bfseries,
    colback=boxBackground, colframe=boxFrame,
    coltitle=black, colbacktitle=boxTitle,
    attach boxed title to top left={xshift=0.3cm,
                                    yshift*=-\tcboxedtitleheight/2},
    boxed title style={
      before upper=\hspace*{0.5cm}, % reserve space for the image
      overlay={
       \node at ([xshift=0.5cm]frame.west)
        
      }
    }
  }
}
\tcbset{my other box/.style={
    enhanced, fonttitle=\bfseries,
    colback=rose-l, colframe=rose-d,
    coltitle=black, colbacktitle=rose,
    attach boxed title to top left={xshift=0.3cm,
                                    yshift*=-\tcboxedtitleheight/2},
    boxed title style={
      before upper=\hspace*{0.5cm}, % reserve space for the image
      overlay={
       \node at ([xshift=0.5cm]frame.west)
        
      }
    }
  }
}

\newtcolorbox{mybox}[1][]{my box, #1}
\newtcolorbox{mechanism}[1][]{my other box, #1}

% Useful packages
\usepackage{amsmath}
\usepackage{graphicx}
\usepackage[colorlinks=true, allcolors=blue]{hyperref}
\usepackage{subcaption}
\usepackage{graphicx}
\title{BW7 - Regulation der Genexpression in Eukaryoten und Spleißen}
\author{Lil'Jana}



\begin{document}
\maketitle
\section{Arten von RNA}
\begin{table}[H]
    \centering
    \begin{tabular}{c c c c}
        hnRNA & heterogene nukleäre RNA & Primärtranskripte &  \\
        mRNA &  Messenger-RNA & proteincodierende RNA & \\
        rRNA & ribosomale RNA & u.a. Proteinbiosynthese & \\
        tRNA &  Transfer-RNA & transportiert Aminosäuren zu Ribosomen &  \\
        snRNA & Small Nuclear RNA &  Bestandteile des Speißosoms  & \\
        (pre-)miRNA & Mikro-RNA & regulatorische RNAs  & \\
        siRNA & Small Interfering RNA & regulatorische RNAs  &  \\
    \end{tabular}
    \caption{Caption}
    \label{tab:my_label}
\end{table}

\section{rDNA}
rDNA stellt eine direkte und spezifische Form der Genexpression dar, bei der die Endprodukte der Genexpression nicht in Proteine, sondern in RNA-Moleküle umgesetzt werden. Diese rRNA-Moleküle sind direkt funktional und tragen zur Bildung der Ribosomen bei, die wiederum für die Translation anderer mRNAs in Proteine verantwortlich sind.

\section{Promotoren}








\newpage
Bitte gibt jeweils eine kurze Definition von den Begriffen: 
Autosomen
\end{document}